% !TeX program = xelatex
% ============================================================
% report/report.tex — Отчёт о курсовом проекте 1 курса магистратуры
% «Системная и программная инженерия» ФКН НИУ ВШЭ.
%
% Структура — «Методические рекомендации по подготовке КП» ОП СПИ:
%   Титульный лист (Прил. 4) → Реферат → Содержание → Определения,
%   обозначения и сокращения → Введение → Основная часть → Заключение →
%   Список использованных источников → Приложения.
% Основная часть зависит от формата работы (выбран при создании проекта):
%   проектно-исследовательский (§5.1) — анализ и постановка, проектирование,
%     реализация и испытания; техдокументация — приложением или ссылкой на
%     репозиторий (§5.2);
%   академический (§5.6) — 3–5 глав: обзорно-постановочная, теоретические
%     результаты, реализация и/или эксперимент, анализ результатов.
% Семинар наставника: первая глава содержит постановку задачи и все разделы
% заявки (цель, задачи, методы, инструменты), реферат — по ГОСТ 7.9/7.32,
% а не аннотация.
%
% Пишите свой текст вместо жёлтых \hseFill{…}; данные о себе, руководителе
% и годе задайте в настройках проекта.
% ============================================================
\documentclass[12pt,a4paper]{extreport}
% !TeX program = xelatex
% ============================================================
% report/preamble.tex — преамбула Отчёта о курсовом проекте магистратуры
% СПИ ФКН НИУ ВШЭ.
%
% Методичка ОП СПИ §5.1/§5.6: структура отчёта о НИР по ГОСТ 7.32 —
% поля 30/15/20/20, Times New Roman 12 пт, интервал 1,5, абзац 1,25 см,
% сквозная нумерация рисунков и таблиц. Чистый документ без штампа и
% рамки. Язык разделов — язык работы (project.lang).
% ============================================================
\newcommand{\hseDefaultMono}{Consolas}
\input{../common/fonts}

((* if project.lang == 'en' *))
\usepackage[russian,english]{babel}
\addto\captionsenglish{\renewcommand{\bibname}{References}}
((* else *))
\usepackage[english,russian]{babel}
\addto\captionsrussian{\renewcommand{\bibname}{Список использованных источников}}
((* endif *))


% Поля страницы — по настройке проекта report_margins (Настройки → Поля
% страницы отчёта). По умолчанию ГОСТ 7.32-2017 (30/15/20/20); вариант ВШЭ —
% 25/10/20/20. Значение подставляется при создании проекта.
((* if 'ВШЭ' in (project.meta.report_margins | default('')) *))
\usepackage[a4paper,left=25mm,right=10mm,top=20mm,bottom=20mm]{geometry}
((* else *))
\usepackage[a4paper,left=30mm,right=15mm,top=20mm,bottom=20mm]{geometry}
((* endif *))
\usepackage{setspace}
\onehalfspacing
\usepackage{indentfirst}
\setlength{\parindent}{1.25cm}
\setlength{\parskip}{6pt}

\usepackage{microtype}
\usepackage{fvextra}
\usepackage{csquotes}
\usepackage{amsmath}
\usepackage{amssymb}
\usepackage{etoolbox}
\usepackage{graphicx}
\usepackage{letltxmacro}
\LetLtxMacro{\oldincludegraphics}{\includegraphics}
\renewcommand{\includegraphics}[2][]{%
  \IfFileExists{#2}{%
    \oldincludegraphics[#1]{#2}%
  }{%
    \IfFileExists{#2.png}{\oldincludegraphics[#1]{#2.png}}{%
      \IfFileExists{#2.jpg}{\oldincludegraphics[#1]{#2.jpg}}{%
        \fcolorbox{red}{white}{%
          \begin{minipage}{0.5\textwidth}\centering
            \textcolor{red}{\textbf{Файл не найден:}} \\ \texttt{\detokenize{#2}}
          \end{minipage}%
        }%
      }%
    }%
  }%
}
\usepackage{xcolor}
\usepackage{hyperref}
\usepackage{lastpage}
\usepackage{float}
\usepackage{array}
\usepackage{booktabs}
\usepackage{longtable}
\usepackage{tabularx}
\usepackage{adjustbox}
\usepackage{xurl}
\urlstyle{same}
\usepackage{hyphenat}
\usepackage{caption}
% Ссылки [n] печатаются в строку (ГОСТ 7.0.5-2008), не надстрочно.
\usepackage{cite}
\DeclareCaptionLabelSeparator{emdash}{~---~}
\captionsetup[table]{position=top,labelsep=emdash,justification=raggedright,singlelinecheck=false}
\captionsetup[figure]{position=bottom,labelsep=emdash}

% ── Листинги кода ───────────────────────────────────────────
\usepackage{listings}
((* if project.lang != 'en' *))\renewcommand{\lstlistingname}{Листинг}((* endif *))
\definecolor{codegray}{rgb}{0.5,0.5,0.5}
\definecolor{backcolour}{rgb}{0.95,0.95,0.92}
\lstset{
  basicstyle=\footnotesize\ttfamily,
  backgroundcolor=\color{backcolour},
  commentstyle=\color{codegray}\itshape,
  numberstyle=\tiny\color{codegray},
  breaklines=true, captionpos=b, numbers=left, numbersep=5pt,
  keepspaces=true, showspaces=false, showstringspaces=false, showtabs=false,
  tabsize=2, frame=single, rulecolor=\color{gray!30},
  xleftmargin=10pt, xrightmargin=10pt,
}

\usepackage{enumitem}
\setlist{nosep}
\setlist[itemize,1]{label=--}
\usepackage{chngcntr}
\counterwithout{table}{chapter}
\counterwithout{figure}{chapter}
\renewcommand{\thetable}{\arabic{table}}
\renewcommand{\thefigure}{\arabic{figure}}
\usepackage{fancyhdr}
\pagestyle{plain}

\usepackage{titlesec}
\usepackage{tikz}
\usetikzlibrary{shapes.geometric, arrows, positioning, calc, er, fit}

\emergencystretch=5em
\tolerance=2000
\hyphenpenalty=300
\exhyphenpenalty=300
\titleformat{\chapter}{\bfseries\centering}{\thechapter}{1em}{}
\titleformat{\section}{\bfseries}{\thesection}{1em}{}
\titlespacing*{\section}{1.25cm}{*2}{6pt}
\titleformat{\subsection}{\bfseries}{\thesubsection}{1em}{}
\titlespacing*{\subsection}{1.25cm}{*2}{6pt}
\titleformat{\subsubsection}{\bfseries}{\thesubsubsection}{1em}{}


% babel переустанавливает \contentsname в \begin{document} — задаём через captions.
((* if project.lang == 'en' *))\addto\captionsenglish{\renewcommand{\contentsname}{Contents}}((* else *))\addto\captionsrussian{\renewcommand{\contentsname}{СОДЕРЖАНИЕ}}((* endif *))

\setcounter{tocdepth}{3}
\setcounter{secnumdepth}{3}

\hypersetup{unicode=true,breaklinks=true,bookmarksopen=true,bookmarksnumbered=true,colorlinks=true,linkcolor=black,urlcolor=blue,citecolor=blue}

% Значки для матриц сравнения (обзор аналогов/подходов).
\newcommand{\cmpplus}{\textbf{+}}
\newcommand{\cmpminus}{\textbf{--}}
\newcommand{\cmpplanned}{\textcolor{blue}{$\pm$}}

% Формат номеров в списке источников с точкой.
\makeatletter
\renewcommand\@biblabel[1]{#1.}
\makeatother

% ── Данные проекта: \hseFill, \signdateblank и все \hse-макросы ──
\input{../common/meta}

\begin{document}
\input{title}

% ============================================================
\chapter*{РЕФЕРАТ}
\addcontentsline{toc}{chapter}{Реферат}
% ============================================================
% ГОСТ 7.32-2017 §5.3, ГОСТ 7.9: сведения об объёме, ключевые слова,
% текст реферата. Реферат передаёт содержание и результаты, а не только
% «о чём» работа (это аннотация).
Отчёт \pageref{LastPage}~с., \ldots~рис., \ldots~табл., \ldots~источн., \ldots~прил.

\medskip
% Ключевые слова: 5–15 слов или словосочетаний, ПРОПИСНЫМИ, через запятую.
\hseFill{КЛЮЧЕВОЕ СЛОВО, КЛЮЧЕВОЕ СЛОВО, КЛЮЧЕВОЕ СЛОВО, КЛЮЧЕВОЕ СЛОВО, КЛЮЧЕВОЕ СЛОВО}

\medskip
Объект исследования --- \hseFill{объект}. Цель работы --- \hseFill{цель}.

В процессе работы \hseFill{что сделано: методы, этапы, эксперименты}.

В результате \hseFill{полученные результаты с цифрами: метрики, характеристики решения}.

Область применения --- \hseFill{где применимы результаты}. \hseFill{степень внедрения, рекомендации, развитие}.

\clearpage
\tableofcontents
\clearpage

% ============================================================
\chapter*{ОПРЕДЕЛЕНИЯ, ОБОЗНАЧЕНИЯ И СОКРАЩЕНИЯ}
\addcontentsline{toc}{chapter}{Определения, обозначения и сокращения}
% ============================================================
В настоящем отчёте применяют следующие термины с соответствующими
определениями и сокращения.

\begin{description}[leftmargin=0pt, labelindent=0pt, font=\bfseries, style=nextline]
  \item[\hseFill{Термин}] \hseFill{определение}.
  \item[\hseFill{АББР}] --- \hseFill{расшифровка сокращения}.
\end{description}

\clearpage
% ============================================================
\chapter*{ВВЕДЕНИЕ}
\addcontentsline{toc}{chapter}{Введение}
% ============================================================
% Актуальность и мотивация, исследовательский вопрос, объект и предмет,
% цель, задачи, методы, практическая значимость, структура отчёта —
% в развитие утверждённой заявки на проект.
\hseFill{актуальность: какая проблема решается, кого она затрагивает и почему существующих решений недостаточно}~\cite{gost732}.

\textbf{Объект исследования} --- \hseFill{объект}.
\textbf{Предмет исследования} --- \hseFill{предмет}.

\textbf{Цель работы} --- \hseFill{цель из заявки на проект}.

Для достижения цели решаются задачи:
\begin{enumerate}[label=\arabic*), leftmargin=1.25cm]
  \item \hseFill{задача 1};
  \item \hseFill{задача 2};
  \item \hseFill{задача 3}.
\end{enumerate}

((* if project.kind == 'research' *))
\textbf{Научная новизна} --- \hseFill{что нового получено по сравнению с известными работами}.
((* else *))
\textbf{Новизна и практическая значимость} --- \hseFill{чем решение отличается от аналогов и где применяется}.
((* endif *))

Отчёт состоит из введения, \hseFill{N} глав, заключения, списка использованных источников и приложений. \hseFill{кратко: что в каждой главе}.

((* if project.kind == 'research' *))
\clearpage
% ============================================================
%  АКАДЕМИЧЕСКИЙ ФОРМАТ (методичка §5.6): 3–5 глав
% ============================================================
\chapter{ОБЗОР ПРЕДМЕТНОЙ ОБЛАСТИ И ПОСТАНОВКА ЗАДАЧИ}
\section{Состояние исследований по теме}
Дайте аналитический обзор работ по теме со ссылками на источники.

\section{Сравнительный анализ подходов}
Сопоставьте известные подходы по значимым критериям (таблица~\ref{tab:approaches}).

\begin{table}[H]
\caption{Сравнение существующих подходов}
\label{tab:approaches}
\centering
\begin{tabular}{|>{\raggedright\arraybackslash}p{5cm}|c|c|c|}
\hline
\textbf{Критерий} & \textbf{Подход~А} & \textbf{Подход~Б} & \textbf{Предлагаемый} \\
\hline
\hseFill{критерий 1} & \cmpplus & \cmpminus & \cmpplus \\
\hline
\hseFill{критерий 2} & \cmpminus & \cmpplanned & \cmpplus \\
\hline
\end{tabular}
\end{table}

\section{Постановка задачи}
Сформулируйте исследовательский вопрос, гипотезу и формальную постановку:
входные данные, ограничения, критерий качества.

\section{Выводы по главе}
\hseFill{нерешённая проблема и постановка задачи исследования}.

\clearpage
\chapter{ОСНОВНЫЕ ТЕОРЕТИЧЕСКИЕ РЕЗУЛЬТАТЫ}
\section{Модель (метод, алгоритм)}
Опишите предложенную модель или метод: допущения, шаги, свойства.

\section{Обоснование}
Приведите доказательства свойств или оценки сложности.

\section{Выводы по главе}
\hseFill{ключевые результаты главы}.

\clearpage
\chapter{РЕАЛИЗАЦИЯ И ЭКСПЕРИМЕНТ}
\section{Программная реализация}
Опишите прототип: архитектуру, технологии, ограничения. Программа --- обязательная
часть исследовательской работы (Программа практики ОП СПИ, §2.1.2).

\section{Методика эксперимента}
План эксперимента, данные (объём, источник, подготовка), метрики.

\section{Выводы по главе}
\hseFill{итоги реализации и эксперимента}.

\clearpage
\chapter{АНАЛИЗ ПОЛУЧЕННЫХ РЕЗУЛЬТАТОВ}
\section{Результаты}
Приведите результаты (таблица~\ref{tab:results}) и сравните с базовыми методами.

\begin{table}[H]
\caption{Результаты эксперимента}
\label{tab:results}
\centering
\begin{tabular}{|>{\raggedright\arraybackslash}p{5cm}|c|c|}
\hline
\textbf{Метод} & \textbf{Метрика~1} & \textbf{Метрика~2} \\
\hline
Базовый метод      & \ldots & \ldots \\
\hline
Предлагаемый метод & \ldots & \ldots \\
\hline
\end{tabular}
\end{table}

\section{Обсуждение и ограничения}
Оцените значимость различий, угрозы достоверности и границы применимости.

\section{Выводы по главе}
\hseFill{ответ на исследовательский вопрос}.
((* else *))
\clearpage
% ============================================================
%  ПРОЕКТНО-ИССЛЕДОВАТЕЛЬСКИЙ ФОРМАТ (методичка §5.1)
% ============================================================
\chapter{АНАЛИЗ ПРЕДМЕТНОЙ ОБЛАСТИ И ПОСТАНОВКА ЗАДАЧИ}
\section{Описание предметной области}
Опишите предметную область, пользователей и их задачи.

\section{Обзор и сравнительный анализ существующих решений}
Сравните аналоги по функциональности, эффективности и удобству
(таблица~\ref{tab:analogues}).

\begin{table}[H]
\caption{Сравнение существующих решений}
\label{tab:analogues}
\centering
\begin{tabular}{|>{\raggedright\arraybackslash}p{5cm}|c|c|c|}
\hline
\textbf{Критерий} & \textbf{Аналог~А} & \textbf{Аналог~Б} & \textbf{Разрабатываемое решение} \\
\hline
\hseFill{критерий 1} & \cmpplus & \cmpminus & \cmpplus \\
\hline
\hseFill{критерий 2} & \cmpminus & \cmpplanned & \cmpplus \\
\hline
\end{tabular}
\end{table}

\section{Постановка задачи и требования}
Сформулируйте задачу, функциональные и нефункциональные требования
(подробно --- в техническом задании, приложение~А).

\section{Выводы по главе}
\hseFill{что именно предстоит разработать и почему}.

\clearpage
\chapter{ПРОЕКТИРОВАНИЕ}
\section{Архитектура системы}
Опишите общую архитектуру, компоненты и их взаимодействие; обоснуйте
выбор решений.

\section{Модель данных и интерфейсы}
Опишите модель данных, программные интерфейсы (API) и пользовательский интерфейс.

\section{Выводы по главе}
\hseFill{ключевые проектные решения}.

\clearpage
\chapter{РЕАЛИЗАЦИЯ И ИСПЫТАНИЯ}
\section{Технологии и инструменты}
Языки, фреймворки, инфраструктура, методология разработки и системы
управления проектом.

\section{Особенности реализации}
Ключевые решения, объёмные характеристики (компоненты, строки кода, тесты).

\section{Испытания и оценка}
Методика испытаний, результаты тестирования, сравнение с аналогами по
выбранным метрикам.

\section{Выводы по главе}
\hseFill{итоги реализации и испытаний}.
((* endif *))

\clearpage
% ============================================================
\chapter*{ЗАКЛЮЧЕНИЕ}
\addcontentsline{toc}{chapter}{Заключение}
% ============================================================
% Выводы по каждой задаче, достижение цели, дальнейшее развитие (ВКР).
В ходе работы:
\begin{enumerate}[label=\arabic*), leftmargin=1.25cm]
  \item \hseFill{итог по задаче 1};
  \item \hseFill{итог по задаче 2}.
\end{enumerate}
Цель работы достигнута: \hseFill{как подтверждено достижение ожидаемых результатов}.
Направления дальнейшей работы: \hseFill{развитие в ВКР}.

\clearpage
% ── Список использованных источников (ГОСТ Р 7.0.5-2008): каждый источник
%    упомянут в тексте. ──
\phantomsection
\addcontentsline{toc}{chapter}{Список использованных источников}
\begin{thebibliography}{99}
  \bibitem{gost732}
  ГОСТ 7.32--2017. Отчёт о~научно-исследовательской работе. Структура и~правила оформления. --- М.: Стандартинформ, 2017. --- 27~с.
\end{thebibliography}

\clearpage
% ── Приложения (ГОСТ 7.32-2017 §6.14), рус. буквами А, Б, В… ──
\appendix
\renewcommand{\thechapter}{\Asbuk{chapter}}
\titleformat{\chapter}[hang]
  {\normalfont\large\bfseries\centering}{Приложение~\thechapter}{1ex}{}
((* if project.kind == 'research' *))
\chapter{Исходный код и данные}
Ссылка на репозиторий с программой и данными эксперимента: \hseFill{URL репозитория в GitLab ФКН}.
((* else *))
\chapter{Техническая документация}
% Методичка §5.2: техническая документация (ТЗ, руководство оператора, ПМИ,
% текст программы, ТЭО/бизнес-модель) приложением ИЛИ ссылкой на систему
% управления проектом / git-репозиторий с планом и документами (wiki, MD).
Техническая документация и план работ ведутся в репозитории проекта:
\hseFill{URL репозитория в GitLab ФКН}. Режим доступа: \hseFill{как получить доступ}.
((* endif *))

\end{document}
